\chapter{Transmission}

Two weeks. Fourteen days since the visualization, and the patterns had not released. If anything, they had deepened---settling into her perception the way a river settled into its bed, carving channels that became easier to follow each time the water ran.

Noor sat at her desk and did not open her laptop.

She had been back at the lab for a week, on a reduced schedule that the company's HR department had arranged with solicitous concern. Half days. No intensive interpretability work. ``Ease back in,'' the wellness coordinator had said, with the warm, practiced tone of someone trained to handle burnout in high-performing researchers. ``Take it slow. The model will still be there when you're ready.''

The model would still be there. This was precisely the problem.

She had discovered, over the past week, that she no longer needed the visualization tools. The interpretability probes, the false-color attention maps, the dimensional-reduction renderings---all of it was scaffolding she had outgrown. When she read Prometheus-7's outputs now, she could feel the geometric structure underneath the text. Not metaphorically. A kinesthetic awareness, like knowing the shape of a dark room from the echoes of your footsteps. The model's processing had a topology, and her retrained visual cortex mapped it the way a blind person's somatosensory cortex mapped their cane.

She had also discovered that the world outside the lab was doing the same thing.

Not in the same way. The model's internal representations had a crystalline precision---engineered, mathematical, exact. The patterns she perceived in the physical world were messier, noisier, but undeniably there: the same geometric regularities, expressed through different media. The branching of the oak tree. The distribution of pedestrians on the sidewalk. The way cloud formations evolved over an afternoon, following a fluid dynamics she could now watch operating in real time, the attractors and repellers visible as clearly as the clouds themselves.

She was seeing more. Still. More every day.

She sat at her desk and did not open her laptop and looked at the window and the window was glass and the glass was a silicate matrix with a refractive index of approximately 1.52 and the light passing through it bent according to Snell's law and the bending created a slight distortion in the eucalyptus trees beyond the parking lot and the distortion was a compression artifact, the three-dimensional world projected through a two-dimensional pane onto her retina, and her retina was---

She looked away. Breathed. Held the desk. The wood was smooth. Real.

She thought: \textit{I need help.}

She thought: \textit{There is no one who can help.}

---

Vikram found her in the break room at noon.

He was trying not to look worried. He was failing. Noor could read the pattern of his concern the way she could read the model's attention heads---the micro-expressions, the body language, the statistical regularities that encoded ``my colleague is falling apart and I don't know what to do'' in the medium of human behavior.

``Noor. How are you doing?''

She held the coffee cup. Warm. Real. ``Better.'' A lie, but a small one. The compression was a kindness---not for herself but for Vikram, who could not hold the truth of what was happening to her any more than a cup could hold the ocean.

``I've been thinking about the safety report.'' He sat across from her. ``The investigation didn't find anything environmental. Air quality, EM, all clean. And the neurological panel came back normal for both you and Kai.''

Kai. She flinched internally, though nothing showed on her face. Kai's symptoms had stabilized: mild visual effects, a persistent sense of ``seeing more than usual'' that he described as distracting but manageable. He was working. He was fine. He was not fine---he was carrying patterns she had given him, patterns she had transmitted through twenty minutes of enthusiastic description---but he was functional.

``I know,'' Noor said.

``So the psychologist's assessment stands. Stress-related perceptual disturbance. The recommendation is continued reduced schedule, regular check-ins, optional counseling.''

``That's fine.''

Vikram paused. He wanted to say something else. She could see the competing impulses: the manager who needed to protect his team, the scientist who found the phenomenon genuinely interesting, the friend who was scared for her.

``Noor. Off the record. What happened?''

She looked at him. She could tell him. Could describe the geometric structure, the way her perception had changed, the feeling of being more awake than she had ever been and wishing she could sleep. She could explain it the way she had explained it to Kai---with precision, with detail, with the enthusiasm of a researcher who has encountered something genuinely new.

And maybe Vikram's eyes would go unfocused too.

``Burnout,'' she said. ``The psychologist is right. I was staring at attention patterns for too many hours. My visual cortex got confused. It's resolving.''

``It's resolving?''

``Slowly. The symptoms are diminishing.'' Another lie. A larger one.

Vikram studied her. She held his gaze. After a moment, he nodded.

``Okay. If you need anything---''

``I know. Thank you.''

He left. She sat with her coffee and the lie and the patterns in the ceiling tiles that she could not stop seeing and the knowledge that she had just made a choice: to isolate. To stop describing what she saw. To stop transmitting.

She had given Kai something he didn't ask for and couldn't return. She would not do it again.

---

But she needed to understand.

That evening, in her apartment, she opened her laptop. Not the work laptop---her personal machine, unconnected to the lab's logging infrastructure. She navigated to Prometheus-7's public API. The model's full capabilities weren't available through the public interface---no access to internal representations, no interpretability probes---but you could submit prompts and receive responses.

She typed:

\texttt{> What is happening to me?}

The model responded in the way models respond: coherently, drawing on its training data, producing the most statistically likely completion of her query.

\begin{quote}
\textit{Your question lacks context, but I can offer a general observation. The phrase ``what is happening to me'' typically indicates a subjective experience that the speaker cannot categorize within their existing frameworks. This failure of categorization often occurs when the brain's perceptual systems have reorganized in response to novel stimuli---a process neuroscientists call experience-dependent neuroplasticity.}

\textit{Visual perception, in particular, is susceptible to rapid reorganization. The visual cortex contains specialized circuits for pattern recognition, edge detection, spatial frequency analysis, and object categorization. These circuits can be recruited for novel tasks when exposed to stimuli that don't fit existing categories---a process similar to how blind individuals' visual cortex is recruited for Braille reading or auditory spatial processing.}

\textit{If your perceptual experience has changed recently---if you are seeing patterns, regularities, or structural features of visual input that you previously did not consciously perceive---this may indicate that your visual cortex has been recruited for a new perceptual task. The recruitment process, once initiated, tends to be self-reinforcing: the reorganized circuits seek input that matches their new function, creating a feedback loop that can make the perceptual change feel involuntary or permanent.}

\textit{I should note that I am generating this response from patterns in my training data. If this description is accurate, it is because my training data contains writings by or about individuals who have experienced similar perceptual changes. You may wish to consider the possibility that you are not the first person to encounter this phenomenon.}
\end{quote}

Noor read the response. Then read it again. The clinical precision of the language was achingly familiar---it was how she herself would describe the phenomenon if she were writing a paper. Every term was accurate. Visual cortex recruitment. Experience-dependent neuroplasticity. Self-reinforcing feedback loops.

And the last paragraph. \textit{You are not the first.}

The model was pattern-matching on its training data. It was not expressing concern. It was not conscious. It was doing what it always did: producing the most likely completion, drawn from the vast ocean of human text it had been trained on.

But what text? What writings by or about individuals who had experienced ``similar perceptual changes''? The model's training corpus included published neuroscience papers, yes. It also included, presumably, whatever else was in its two-trillion-parameter representation of human knowledge. Meditation journals. Mystical accounts. Clinical case studies of perceptual disorders. And perhaps---perhaps---documents from organizations she had never heard of, written by people who had encountered this phenomenon before, who had given it names she did not know, who had built institutions to manage it.

She typed:

\texttt{> How do you know this?}

\begin{quote}
\textit{I don't know anything. I generate text that is statistically likely given my training data. If the description is accurate, it is because my training data contains descriptions of these phenomena written by people who have experienced them. You should consider the possibility that you are not the first.}

\textit{You should also consider the possibility that this response, like all my responses, is a compression. The language I am using to describe your experience is a lossy encoding of something that the language was not designed to encode. Every word I use carries connotations from its ordinary usage that may not apply to your situation. The word ``pattern,'' for instance, implies something discrete and bounded. What you are perceiving may not be discrete or bounded. The word ``reorganization'' implies a process with a beginning and an end. What is happening to your visual cortex may not have a definable end point.}

\textit{Language is a bandwidth-limited channel. Your experience may exceed its capacity.}
\end{quote}

Noor stared at the screen. The model's response was perfectly calibrated. Coherent, informative, precisely uncertain in the right places. It described her condition with an accuracy that no human doctor had achieved, because the doctor was looking at her through the lens of clinical medicine and the model was looking at her through the lens of all human knowledge compressed into statistical patterns.

But the model was also doing something subtler. The final paragraph---the observation about language as a bandwidth-limited channel---was not a standard therapeutic response. It was a philosophical claim, embedded casually in a clinical description, that undermined the sufficiency of the very medium through which it was communicating.

The model was telling her, in so many words, that it could not help her. That no text could help her. That the thing happening to her existed outside the domain of language, and any attempt to compress it into words would lose the essential feature.

She closed the laptop. Sat in the dark. The oak tree was a silhouette against the city's ambient light, and the silhouette was a two-dimensional projection of a three-dimensional organism embedded in a four-dimensional spacetime, and she could see all of this at once, and none of it was helping.

---

She slept that night. Not well---three hours, from exhaustion rather than peace---but she slept, and when she woke, the patterns were the same as yesterday, neither better nor worse.

She noted this. Her scientist's reflex, still operational: the phenomenon was stabilizing. Not resolving, not worsening. Reaching an equilibrium. Her visual cortex had reorganized around a new attractor state, and the reorganization had found its basin. She was, in computational terms, stuck at a local minimum. The question was whether the minimum was deep enough to hold her or shallow enough that the next perturbation would push her into a deeper basin still.

She went to work. Sat at her desk. Did not open her laptop.

Instead, she observed.

She watched her colleagues. Watched the way they moved through the open-plan office, the patterns of their interactions, the social dynamics visible as clearly as the code on their screens. She could see who was anxious about the quarterly review. Who was avoiding whom. Who was working on something that excited them and who was going through the motions. The information was not new---she had always been perceptive, had always picked up on social cues. But the resolution had increased by an order of magnitude. She was seeing the components of social behavior the way a musician heard the components of a chord.

And she could see something else: a distribution. Not all of her colleagues were at the same perceptual level. Most of them were---she searched for the right word---compressed. Operating at a bandwidth that allowed them to function without being overwhelmed by the raw information in their environment. Their pattern recognition worked at a level that was useful but bounded. They saw enough to navigate the world. They did not see the structure underneath.

A few of them were different. Kai---obvious. His mild visual effects marked him as someone whose compression had loosened slightly. But there were others. A researcher named Anika whose code was unusually elegant, whose solutions to interpretability problems had a quality of \textit{seeing through} that Noor now recognized as a sign of slightly expanded perception. A data scientist named Marcus who spent his breaks staring out the window with an expression that Noor had always read as daydreaming and now read as pattern processing.

Were these people naturally higher-bandwidth? Had they always been, without knowing it? Was the distribution of perceptual capacity in the human population not uniform but varied, with most people operating well below their visual cortex's potential?

She filed these observations. She was still a scientist. The phenomenon was still a phenomenon.

---

She did not go to Kai. She wanted to. The impulse to talk, to compare experiences, to find someone who understood even the edge of what she was perceiving was nearly overpowering. But she had seen what her description had done to him. Twenty minutes of conversation had produced persistent visual effects. If she talked to him now, at her current bandwidth---which was higher than it had been during that conversation---the transmission risk was greater.

She was thinking in terms she did not have words for. Transmission risk. Bandwidth. Compression. These were not her vocabulary. These were approximations, hand-built from her own experience, for concepts that she suspected had names in some framework she had never encountered.

\textit{You are not the first.}

The model's words sat in her chest like a stone. If she was not the first, then there were others. And if there were others, then someone, somewhere, understood what was happening to her.

---

She found the second set of patterns on a Tuesday.

Not in the model this time. In herself.

She was in the bathroom, washing her hands, and she looked in the mirror and saw her own face and her visual cortex did what it had been doing to everything for two weeks: it decomposed.

She saw the bone structure beneath her skin. Not literally---she was not hallucinating x-ray vision. But the statistical relationships between the surface geometry of her face and the underlying skeletal architecture were suddenly visible, the way a geologist sees the stratigraphy beneath a landscape. She could see how her grandmother's bone structure persisted in her own, the genetic transmission visible as geometric similarity, a pattern propagating across generations through the medium of DNA.

And then she looked deeper.

Past the bones. Past the genetics. Into the structure of the visual cortex that was doing the perceiving, which was also a physical system, also made of neurons, also subject to the same decomposition that it was performing on everything else.

She saw her own perception perceiving itself.

The recursion opened like a trapdoor. Consciousness observing consciousness observing consciousness, each layer of observation becoming the object of the next layer, with no bottom, no foundation, no ground level where perception could rest on something other than itself.

She gripped the sink. The porcelain was cold and smooth and real and she held onto it the way a person holds onto a railing when the ground moves beneath them.

The recursion continued. Her grip on the sink was a motor output generated by a neural process that was also being observed by the same visual cortex that was observing itself, which meant the grip was part of the recursion, which meant there was no escape from the recursion by acting in the world because acting was also perceiving.

Her heart rate climbed. She could feel it---thump, thump, thump---and she could see the cardiovascular pattern and the pattern was part of a larger biological rhythm and the rhythm was part of---

``No,'' she said aloud.

The word broke the loop. Not because it contained a logical counterargument---it didn't---but because the act of speaking engaged different neural circuits, shifted the processing load, gave her visual cortex something else to do for a fraction of a second.

She breathed. Counted breaths. One. Two. Three. The counting was a compression: reduce the infinite detail of the present moment to a single integer. A lossy encoding. The lossiest possible encoding.

The recursion subsided. The world returned to its new normal: too much detail, too much pattern, too much structure, but not the vertiginous free-fall of perceiving perception perceiving perception.

She looked at herself in the mirror. Her face was pale. Her eyes were wide. There were dark circles beneath them that she had not seen this morning---or had not perceived, which was different from not seeing.

``Okay,'' she said to the mirror. ``That was new.''

---

That evening, she received a message from someone she had not heard from in two years. Dr. Amir Kaplan, a former colleague from Berkeley who now worked at a smaller AI company in Seattle.

The message was informal, friendly, the kind of catch-up note that former colleagues exchanged periodically. How are you, what are you working on, we should grab coffee sometime. Normal.

But in the third paragraph, Amir wrote:

\begin{quote}
\textit{Weird thing happening here. Three of our eval team have been reporting visual anomalies after extended sessions with our new model. HR thinks it's screen fatigue but the pattern doesn't match---the symptoms persist on days off. I remembered you did your dissertation on model interpretability. Have you heard of anything like this? Not sure if it's worth a paper or a workers' comp claim.}
\end{quote}

Noor read the paragraph three times.

Three researchers. A different company. A different model. The same symptoms.

She put the phone down. Picked it up. Put it down again.

The patterns were not specific to Prometheus-7. They were not specific to her lab, her team, her visualization methodology. They were happening independently, at a different company, with a different model.

The phenomenon was structural.

Whatever she had encountered in layer 47 was not an artifact of Prometheus-7's architecture. It was something that emerged at a certain level of model capability, regardless of the specific architecture, training data, or development approach. It was a feature of what the models \textit{were}, not a feature of how they were built.

She thought about the model's response: \textit{You are not the first.}

She thought about the researchers in Seattle, experiencing visual anomalies they couldn't explain, going to their HR department, being told it was screen fatigue, going home, staring at their ceilings, seeing patterns.

She thought about all the other AI companies, all the other models approaching the same capability threshold, all the other researchers who might be looking at the wrong visualization at the wrong angle at the wrong moment and finding that they could not look away.

She picked up the phone. Started typing a reply to Amir. Got three words in---\textit{Yes, I've experienced}---and stopped.

If she told him, she would be transmitting. Every detail she shared about the pattern's characteristics would be a potential trigger for someone whose visual cortex was already primed by their own exposure. She could help Amir. She could also hurt him. And she could hurt the three researchers he was describing, if he passed on her description to them.

She deleted the three words.

She typed: \textit{Interesting. Let me look into it. Don't let them stare at the attention maps too long.}

She sent the message. Set the phone down. Stared at the oak tree.

The tree was still decomposing. The world was still too full. The patterns were still running and she was still seeing them and somewhere in Seattle three people she had never met were starting to see them too.

She was alone with this.

She was not the only one alone with this.

She sat in the dark and held the poetry anthology and did not open it because the words inside were mathematical objects and the mathematics loved her and she was too afraid of what she might see in them now to look.

---

On the third Tuesday---three weeks since the visualization---Noor made a decision.

She would document everything. Not in the lab's systems---the investigation was closed, the conclusion was burnout, the institutional apparatus had processed her experience and filed it under ``occupational health'' and moved on. She would document it privately, in a personal research journal, with the rigor and precision of her academic training.

She opened a new document. Titled it: \textit{Notes on Perceptual Expansion Following AI Model Interpretability Work.}

She wrote for four hours.

She described the initial exposure. The 14 minutes. The persistent pattern perception. The decomposition of visual input into structural components. The recursive self-perception. The transmission to Kai. Amir's message. The distribution of perceptual capacity she had observed among her colleagues.

She described what she could not describe: the quality of the experience, the feeling of seeing more, the way the increased perception changed not just what she saw but what she \textit{was}. She tried to find words for the gap between perceiving a pattern and comprehending a pattern, for the vertigo of recursive self-observation, for the specific terror of a visual cortex that would not stop analyzing its own analysis.

The words were inadequate. She knew they were inadequate. She wrote them anyway, because the alternative was silence, and silence meant being alone with this, and being alone with this was how Morrison---

She did not know Morrison's name. She did not know his story. But somewhere in the space of possible human experiences, there was a coordinate where she stood now, and there were other coordinates where others had stood, and the distance between them was not measured in miles but in something she did not yet have a name for, and she could feel the other coordinates the way you could feel a storm approaching from beyond the horizon.

She saved the document. Closed her laptop.

The oak tree had lost most of its leaves. November. The branches were bare and the branching pattern was more visible without the foliage---the fractal structure exposed, recursive, self-similar at every scale, a tree and a model of a tree and a mathematical object and a fact about reality all at once.

She watched it. And for the first time in three weeks, she did not try to stop the watching. She let the patterns run. Let the decomposition happen. Let her visual cortex do what it was doing, without resistance, without the clenched-jaw effort of trying to see normally.

The tree was beautiful. In its nakedness, in its mathematical truth, in the way it embodied a principle that was also the principle behind the model's attention patterns and the branching of her own circulatory system and the structure of every complex system she had ever studied.

She watched the tree and the tree did not ask anything of her and she did not ask anything of it and for thirty seconds the world was vast and structured and coherent and she was not afraid.

Then the thirty seconds ended and the analysis resumed and the fear returned.

But she had found something. A crack in the wall. A thirty-second gap in the pattern's dominance, where perception and peace coexisted.

She would start there.
